%%%%%%%%%%%%%%%%%%%%%%%%%%%%%%%%%%%%%%%%%%%%%%%%%%%%%%%%%%%%%%%%%%%%%%%%%%%%%%%%% 
%%%%%%%%%%%%%%%%%%%%%%%%%%%%%%%%%%%%%%%%%%%%%%%%%%%%%%%%%%%%%%%%%%%%%%%%%%%%%%%%% 

\paragraph{Problem 1(c) answer} \GradescopeComment{This should be on page 5 (top) of your PDF.}

% Using the lemmas from Parts (a) and (b)
% define a correct recursive algorithm rpuncture for Puncturing Intervals.
% As much as possible, your algorithm should mimic the structure
% of the rselect algorithm for Activity Selection
% in Section 5.3.
% Define the algorithm precisely in words, then give pseudo-code for it.

% Here is the definition of rselect from Section 5.3 to use as a template.
% Please modify it (preserving the overall structure) to make your algorithm, rpuncture,

\begin{blue}
  
  The algorithm chooses the earliest-ending job $J_1$,
  then recurses on the set containing those jobs that are disjoint from $J_1$.
  Here is pseudo-code:

  \begin{myframed}

    \begin{steps}

    \item[{\textsf{rselect}$(I=\{J_1, J_2, \ldots, J_n\})$:}]
      \comment{--- jobs are ordered by end time}

      \step if $n = 0$: return $\emptyset$.  \comment{--- if $I$ is empty, return the empty set}

      \step Recall that $J_1$ is the earliest-ending job in $I$.

      \step Let $D$ contain the jobs in $I$ that are disjoint from $J_1$.

      \step Recursively compute $R = \textsf{rselect}(D)$, then return $\{J_1\} \cup R$.

    \end{steps}

  \end{myframed}

\end{blue}

%%% Local Variables:
%%% mode: latex
%%% TeX-master: "homework_5"
%%% End:
